
\section{Principes généraux}


\begin{frame}
\frametitle{Vocabulaire}

Rappel: une \alert{partition} d'un ensemble $S$ est un ensemble $\{F_1, F_2, \cdots F_n\}$
de parties de $S$ telles que $F_i \cap F_j = \emptyset$ et $\bigcup_i Fi = S$.

\vfill

Dans notre cas
\begin{redbullet}
   \item l'ensemble est une collection;
   \item les éléments de l'ensemble sont les documents;
   \item les parties sont nommées \alert{fragments} (\textit{shard} en anglais)
\end{redbullet}

\end{frame}

\begin{frame}
\frametitle{Le partitionnement (\textit{sharding})}

Le partitionnement consiste à  

\begin{redbullet}
  \item \alert{découper} une (grande) collection
      en \alert{fragments} en fonction d'une \alert{clé}
      \item placer chaque fragment sur un \alert{serveur}
      \item Maintenir un \alert{répertoire} indiquant  que telle \alert{clé} se trouve dans tel \alert{fragment}
          sur tel \alert{serveur}
\end{redbullet}
 
\vfill

Le partitionnement apporte la \alert{répartition de charge} (\textit{load balancing})

\vfill

Il doit être \alert{dynamique} (ajout/retrait de serveurs) pour s'adapter ``élastiquement'' 
 
 \vfill
 
 S'applique à des collections de paires (clé, valeur), où valeur est n'importe quelle information structurée.
\end{frame}




\begin{frame}
\frametitle{La clé de partitionnement}

La clé indique à quel fragment appartient un document.

\vfill

Si un grand nombre de documents partagent la même clé, ils ne peuvent être que
dans le même fragment $\Rightarrow$ le système perd en souplesse de distribution.

\vfill

La valeur de la clé pour un document ne devrait pas être modifiable (nécessité de déplacement).

\vfill

\alert{L'identifiant séquentiel d'un document est un bon choix}

\end{frame}


\begin{frame}
\frametitle{Opérations}

Les opérations de type ``dictionnaire'' peuvent être transmises à \alert{un seul} serveur\,:

\begin{redbullet}
   \item $get(k): v$, recherche par clé
   \item $put (k, v)$, insertion
   \item $delete(k)$, suppression
   \item $range(k_1, k_2)$, recherche par intervalle (peut impliquer plusieurs serveurs, ne marche pas avec le hachage)
\end{redbullet}

\vfill

\alert{Pour toutes les autres opérations}\,: tous les serveurs effectuent l'opération localement.

\vfill

Exemple\,: recherche sur une clé secondaire, s'effectue sur \alert{tous} les serveurs, avec si possible des index locaux.
\end{frame}


\begin{frame}
\frametitle{Techniques de partitionnement}

\begin{redbullet}
   \item par \alert{intervalle}\,: documents triés sur la clé, puis groupés par intervalles.
   
     Exemples représentatifs\,: \mongodb/, HBase (BigTable).
     
   \item par \alert{hachage} sur la clé, avec répertoire de hachage.
   
       Exemples\,: \textit{memcached}, \textit{chord}, \textit{Dynamo/S3/Voldemort}, beaucoup d'autres.
 \end{redbullet}
 
\vfill

\begin{remblock}{Des méthodes connues dans d'autres contextes}
Se référer aux techniques d'indexation par arbre-B et par hachage dans des environnements centralisés.
\end{remblock}
\end{frame}

\begin{frame}
\frametitle{Structure d'un système partitionné}

\figSlideWithSize{partition-vision}{10cm}
      
\begin{redbullet}
  \item Par intervalle = séquentiel indexé, \alert{Arbre B}, Arbre B+, etc. 
  \item Par hachage = hachage statique, dynamique, \alert{hachage linéaire}.
\end{redbullet}

\vfill

\begin{remblock}{Tout en RAM}
La structure de routage doit tenir en mémoire pour être pleinement efficace.
\end{remblock}
\end{frame}



\section{Et en distribué?}


\begin{frame}
\frametitle{Routage et stockage en distribué}

\figSlideWithSize{partition-distrib}{10cm}
      
\vfill

\begin{redbullet}
  \item Routeur = \textit{Master}, \textit{Balancer}, \textit{Config server}, ...
  \item Fragment = \textit{chunk}, \textit{tablet}, \textit{region}, \textit{bucket}, ...
\end{redbullet}

\vfill

\begin{remblock}{Et la réplication?}
Chaque serveur gère une réplication locale pour tolérance aux pannes.
\end{remblock}
\end{frame}

